% sec-01: two institutions, seven capability criteria, and where functions sit.
% Figure 8 (diagrams-type clustered topology), Table 12.
\section{Two Routes, Run in Parallel and Disclosed}

Two La Jolla institutions are approached on the same day by two different routes,
asking two different questions. Each letter discloses the other. Neither is asked
for exclusivity and neither is described to the other as preferred.

\subsection{Why parallel rather than sequential}

A sequential approach costs a quarter. A first institution takes four to eight
weeks to route an unsolicited feasibility request to the right surgeon, and if
the answer is no, the second institution starts from zero after that. For an
eighteen-participant Phase 1 whose protocol is already complete and deposited,
spending a quarter to learn a routing answer is the expensive option
\cite{onpremwhippl}.

A parallel approach costs something too, and it is worth naming. Two institutions
in one small clinical community will learn of each other's involvement. The only
question is whether they learn it from the letter or from a colleague, and
learning it from the letter is strictly better for everyone.

\begin{apptable}
\begin{tabularx}{\textwidth}{>{\raggedright\arraybackslash}p{3.6cm} >{\raggedright\arraybackslash}p{3.4cm} >{\raggedright\arraybackslash}X}
\headrow \thead{Criterion} & \thead{Route one} & \thead{Route two}\\
\midrule
Entry point & Escalation on a published path, after an August 2026 request & An inbound visibility notice to a digital trials group \\
What leads the letter & The clinical program & The artificial intelligence method \\
The ask & A routing decision: who is the right site principal investigator & A read of the public portfolio and a view on the method \\
Pancreatic surgical volume & The criterion a first meeting establishes & The criterion a first meeting establishes \\
Robotic platform availability & Asked, not assumed. No configuration is asserted & Asked, not assumed \\
Trial operations and budgeting & A trials office and a budget questionnaire route exist & A centralized research operations office exists \\
What neither is asked & To endorse, to host, or to commit & To endorse, to host, or to commit \\
\end{tabularx}
\end{apptable}
\vspace{-0.60cm}
\tabcap{Table 12. The two candidate routes against seven criteria, with the two\\
rows that are asked rather than assumed marked by their own wording.}

The asks differ because the entry points differ. Approaching a digital trials
group with a request to host an experimental surgical operation would be asking
the wrong department for the wrong thing, and approaching a cancer center's
leadership with a method-evaluation invitation would waste an escalation that
already has a published path.

\subsection{Where each function would physically sit}

Figure 8 answers the question a site operations lead asks first, which is not who
is responsible but where the work happens. A function that is legally the
sponsor's but physically on a hospital floor needs space, badge access, and a
network drop, and none of that appears in an obligation diagram.

\begin{appfloat}[!tb]
\begin{appfig}[diagrams-type / clustered infrastructure with vector glyphs]
% Three dashed clusters on a 4.55cm pitch; tiles on a 1.75cm pitch.  \dgnode
% places the tile, its pictogram and its label as three related nodes, so a fit
% over the tile and its label encloses both.
\dgnode{dgtile}{0}{0}{\glyphrobot}{Robotic procedure suite}{c1}
\dgnode{dgtile}{0}{-1.75}{\glyphpill}{Investigational pharmacy}{c2}
\dgnode{dgtile}{0}{-3.50}{\glyphuser}{Consent and follow-up clinic}{c3}
\dgnode{dgtile}{0}{-5.25}{\glyphflask}{Pathology and specimens}{c4}
\dgnodew{dgtilem}{4.55}{0}{\glyphai}{On-premises model host}{s1}
\dgnodew{dgtilem}{4.55}{-1.75}{\glyphgear}{Verification harness}{s2}
\dgnodew{dgtilem}{4.55}{-3.50}{\glyphdoc}{Regulatory binder}{s3}
\dgnodeg{dgtileg}{9.10}{0}{\glyphnet}{Robot control network}{x1}
\dgnodeg{dgtileg}{9.10}{-1.75}{\glyphdb}{Electronic data capture}{x2}
\node[dgctitle,anchor=south west] (t1) at (-1.30,0.62) {Clinical campus, candidate site, feasibility stage only};
\node[dgctitle,anchor=south west] (t2) at (3.25,0.62) {Sponsor premises};
\node[dgctitle2,anchor=south west] (t3) at (7.80,0.62) {Reachable by neither};
\node[dgcluster,fit={(c1)(c4l)(t1)},inner sep=7pt] {};
\node[dgcluster,fit={(s1)(s3l)(t2)},inner sep=7pt] {};
\node[dgcluster2,fit={(x1)(x2l)(t3)},inner sep=7pt] {};
\draw[dgedgeb] (s1) -- (s2);
\draw[dgedge] (c1) -- (c4);
\draw[dgedged] (s1) -- (x1);
\draw[dgedged] (s1) -- (x2);
\pxmark{6.95}{0}
\pxmark{6.95}{-1.05}
\node[font=\tiny\itshape,text=fundgray,anchor=north west,text width=118mm]
  at (-1.30,-6.50) {The third cluster is the point of the figure. The model
  process holds no write credential to data capture and no route to robot
  control, which is a property of the wiring rather than of a policy document.};
\end{appfig}
\vspace{-0.60cm}
\figcaption{Figure 8. Where each trial function would physically sit, and the two systems\\
the advisory model process is architecturally prevented from ever reaching.}
\end{appfloat}

\subsection{The clinical introductions that are deliberately deferred}

Two clinicians at one of the two institutions are relevant to this program by
specialty, and no public professional address was located for either. This
company will not guess one. If that institution responds, the introduction is
requested through its own research services rather than attempted directly. The
rule holds generally: an introduction requested through an institution is an
introduction the institution knows about \cite{scrippsdtc2026}.

\subsection{The escalation rule, applied to both}

If there is no response within three business days, the original request is
forwarded once, with a short note, to the contact one level up on the published
path. Once, and not repeatedly. A second escalation inside a week converts a
routing question into an impression of the company, and the impression outlasts
the question.
